\documentclass{article}
\usepackage{amsmath,amssymb,amsthm}
\usepackage{algorithm}
\usepackage[noend]{algpseudocode}
\usepackage{hyperref}
\usepackage{enumitem}
\usepackage{bm}
\usepackage{geometry}
\geometry{margin=1in}
\newtheorem{theorem}{Theorem}
\newtheorem{lemma}{Lemma}
\newtheorem{assumption}{Assumption}
\newcommand{\E}{\mathbb{E}}
\newcommand{\R}{\mathbb{R}}
\newcommand{\norm}[1]{\left\|#1\right\|}
\newcommand{\grad}{\nabla}
\newcommand{\Id}{\mathbf{I}}

\begin{document}

\section*{Algorithm Name}
\textbf{AdaGrad for Non‑Convex Smooth Optimization (Running‑Sum Variant)}  

The method keeps a cumulative sum of past gradients (or of their elementwise squares) and uses this sum to adapt the learning rate per coordinate.

\vspace{0.5em}
\section*{Mathematical Formulation}
Let $f:\R^{d}\to\R$ be the objective function and denote by
\[
g_{k}:=\grad f(w_{k})\quad\text{(or an unbiased stochastic estimate thereof).}
\]
Define the cumulative sum of gradients
\[
A_{k}:=\sum_{i=1}^{k} g_{i}\in\R^{d},
\]
and the cumulative sum of elementwise squared gradients
\[
v_{k}:=\sum_{i=1}^{k} g_{i}^{\odot 2}\in\R^{d},
\qquad
g_{i}^{\odot 2}:=(g_{i,1}^{2},\dots,g_{i,d}^{2})^{\top}.
\]
Given a base stepsize $\alpha>0$ and a small constant $\varepsilon>0$ for numerical stability, the update at iteration $k\ge 0$ is
\begin{equation}
\label{eq:update}
w_{k+1}=w_{k}-\eta_{k}\odot g_{k},
\qquad 
\eta_{k}:=\frac{\alpha}{\sqrt{v_{k}}+\varepsilon},
\end{equation}
where the division and square root are taken elementwise.  
When $g_{k}$ is stochastic, we write $g_{k}=\grad f(w_{k})+\xi_{k}$ with $\E[\xi_{k}\mid w_{k}]=0$.

\vspace{0.5em}
\section*{Pseudocode}
\begin{algorithm}[H]
\caption{AdaGrad (Running‑Sum Variant)}\label{alg:adagrad}
\begin{algorithmic}[1]
\State \textbf{Input:} initial point $w_{0}\in\R^{d}$, base stepsize $\alpha>0$, stability $\varepsilon>0$, number of iterations $T$.
\State $v_{0}\gets \mathbf{0}\in\R^{d}$ \Comment{initialize accumulated squared gradients}
\For{$k=0,\dots,T-1$}
    \State Sample a stochastic gradient $g_{k}$ (or compute the full gradient) at $w_{k}$.
    \State $v_{k+1}\gets v_{k}+g_{k}^{\odot 2}$.
    \State $\eta_{k}\gets \dfrac{\alpha}{\sqrt{v_{k+1}}+\varepsilon}$ \Comment{elementwise division}
    \State $w_{k+1}\gets w_{k}-\eta_{k}\odot g_{k}$.
\EndFor
\State \textbf{Return} $w_{T}$ (or the average $\frac1T\sum_{k=0}^{T-1}w_{k}$).
\end{algorithmic}
\end{algorithm}

\vspace{0.5em}
\section*{Dry Run Example}
Consider the scalar quadratic $f(w)=(w-3)^{2}$, initial point $w_{0}=0$, base stepsize $\alpha=0.1$, and $\varepsilon=10^{-8}$ (ignored for brevity).  
The exact gradient is $g_{k}=2(w_{k}-3)$.

\begin{center}
\begin{tabular}{c|c|c|c|c}
$k$ & $w_{k}$ & $g_{k}=2(w_{k}-3)$ & $v_{k+1}=v_{k}+g_{k}^{2}$ & $w_{k+1}=w_{k}-\dfrac{\alpha}{\sqrt{v_{k+1}}}\,g_{k}$\\\hline
0 & $0$ & $-6$ & $36$ & $0-\dfrac{0.1}{\sqrt{36}}(-6)=0+0.1\cdot\frac{6}{6}=0.1$\\\hline
1 & $0.1$ & $-5.8$ & $36+33.64=69.64$ & $0.1-\dfrac{0.1}{\sqrt{69.64}}(-5.8)
      =0.1+0.1\cdot\frac{5.8}{8.345}\approx0.1695$\\\hline
2 & $0.1695$ & $-5.661$ & $69.64+32.05\approx101.69$ &
 $0.1695-\dfrac{0.1}{\sqrt{101.69}}(-5.661)
 \approx0.1695+0.1\cdot\frac{5.661}{10.084}\approx0.2301$\\\hline
\end{tabular}
\end{center}
The objective values after each step are $f(w_{1})\approx8.41$, $f(w_{2})\approx8.04$, $f(w_{3})\approx7.71$, showing a gradual decrease.

\vspace{1em}
\section*{Assumptions}
\begin{assumption}[Smoothness]\label{ass:smooth}
$f:\R^{d}\to\R$ is $L$‑smooth, i.e.
\[
\forall x,y\in\R^{d}:\quad
f(y)\le f(x)+\langle\grad f(x),y-x\rangle+\frac{L}{2}\norm{y-x}^{2}.
\]
\end{assumption}
\begin{assumption}[Unbiased Stochastic Gradients]\label{ass:unbiased}
For each iteration $k$, the stochastic gradient $g_{k}$ satisfies
\[
\E[g_{k}\mid w_{k}]=\grad f(w_{k}),\qquad
\E\big[\|g_{k}-\grad f(w_{k})\|^{2}\mid w_{k}\big]\le\sigma^{2}.
\]
\end{assumption}
\begin{assumption}[Bounded Gradients]\label{ass:bounded}
There exists $G>0$ such that $\|g_{k}\|\le G$ almost surely for all $k$.
\end{assumption}
All hyper‑parameters satisfy $\alpha>0$, $\varepsilon>0$, and the iteration budget $T\in\mathbb{N}$.

\vspace{0.5em}
\section*{Main Convergence Theorem}
\begin{theorem}\label{thm:main}
Under Assumptions \ref{ass:smooth}–\ref{ass:bounded}, the iterates generated by Algorithm~\ref{alg:adagrad} satisfy
\[
\frac{1}{T}\sum_{k=0}^{T-1}\E\!\left[\big\|\grad f(w_{k})\big\|^{2}\right]
\;\le\;
\frac{2L\big(f(w_{0})-f^{\star}\big)}{\alpha\sqrt{T}}
\;+\;
\frac{2\alpha\sigma^{2}\sqrt{d}}{\sqrt{T}}
\;+\;
\frac{2\alpha G^{2}\sqrt{d}}{\sqrt{T}},
\]
where $f^{\star}=\inf_{w}f(w)$. Consequently,
\[
\min_{0\le k<T}\E\!\left[\big\|\grad f(w_{k})\big\|^{2}\right]
=O\!\left(\frac{\log T}{\sqrt{T}}\right).
\]
\end{theorem}

\vspace{0.5em}
\section*{Theoretical Properties}
\begin{itemize}[leftmargin=*]
\item \textbf{Stability.} The adaptive denominator $\sqrt{v_{k}}$ is non‑decreasing elementwise, guaranteeing that the effective stepsize $\eta_{k}$ never exceeds $\alpha/\varepsilon$ and diminishes as more gradient information accumulates.
\item \textbf{Hyper‑parameter Sensitivity.} The bound depends linearly on $\alpha$ and on $\sqrt{d}$, reflecting the well‑known dimension‑dependence of AdaGrad. Smaller $\alpha$ reduces the variance term but slows down the deterministic term.
\item \textbf{Comparison to SGD.} For SGD with a fixed stepsize $\eta$, the standard non‑convex bound is $O(1/\sqrt{T})$. AdaGrad enjoys the same asymptotic rate while automatically adapting to the geometry of the data; in sparse settings $v_{k}$ grows slowly, yielding larger effective stepsizes.
\item \textbf{Special Cases.} If $f$ is convex, the same proof yields the classic $O(\frac{1}{\sqrt{T}})$ regret bound. If gradients are deterministic ($\sigma=0$) the bound reduces to $O(1/\sqrt{T})$ without the variance term.
\end{itemize}

\vspace{0.5em}
\section*{Discussion}
The theorem shows that even without convexity, the AdaGrad update (with the running‑sum stepsize) drives the average squared gradient to zero at a rate comparable to the optimal $O(1/\sqrt{T})$ for first‑order methods. The proof hinges on the smoothness inequality, the unbiasedness of stochastic gradients, and a careful telescoping of the adaptive stepsizes. The $\sqrt{d}$ factor is unavoidable for elementwise AdaGrad; alternative matrix‑based adaptations can improve this dependence.

\vspace{1em}
\section*{Preliminaries (Definitions and Lemmas)}
\begin{lemma}[Smoothness Inequality]\label{lem:smooth}
If $f$ is $L$‑smooth then for any $x$ and any vector $u$,
\[
f(x-u)\le f(x)-\langle\grad f(x),u\rangle+\frac{L}{2}\|u\|^{2}.
\]
\end{lemma}
\begin{proof}
Directly from Definition~\ref{ass:smooth} with $y=x-u$.
\end{proof}

\begin{lemma}[Summation of Adaptive Stepsizes]\label{lem:sum_eta}
For any non‑negative sequence $\{a_{k}\}_{k=1}^{T}$ define $s_{k}=\sum_{i=1}^{k}a_{i}$. Then
\[
\sum_{k=1}^{T}\frac{a_{k}}{\sqrt{s_{k}}}\le 2\sqrt{s_{T}}.
\]
\end{lemma}
\begin{proof}
Observe that $\frac{a_{k}}{\sqrt{s_{k}}}=2\big(\sqrt{s_{k}}-\sqrt{s_{k-1}}\big)$ and sum telescopically.
\end{proof}

\vspace{0.5em}
\section*{Main Proof (Step‑by‑Step Derivation)}
We prove Theorem~\ref{thm:main}. All expectations are taken w.r.t. the randomness up to iteration $k$ unless otherwise noted.

\begin{enumerate}[label={[\textbf{Step \arabic*.}]}, leftmargin=*]
\item \textbf{One‑step progress via smoothness.}\\
Using Lemma~\ref{lem:smooth} with $u=\eta_{k}\odot g_{k}$,
\[
f(w_{k+1})\le f(w_{k})-\langle\grad f(w_{k}),\eta_{k}\odot g_{k}\rangle
+\frac{L}{2}\norm{\eta_{k}\odot g_{k}}^{2}. \tag{1}
\]

\item \textbf{Take conditional expectation.}\\
Condition on $\mathcal{F}_{k}=\sigma(w_{0},g_{0},\dots,g_{k-1})$. Using unbiasedness \eqref{ass:unbiased},
\[
\E\!\left[\langle\grad f(w_{k}),\eta_{k}\odot g_{k}\rangle\mid\mathcal{F}_{k}\right]
= \big\langle\grad f(w_{k}),\eta_{k}\odot\grad f(w_{k})\big\rangle.
\tag{2}
\]
Similarly,
\[
\E\!\left[\norm{\eta_{k}\odot g_{k}}^{2}\mid\mathcal{F}_{k}\right]
= \norm{\eta_{k}\odot\grad f(w_{k})}^{2}
+\E\!\left[\norm{\eta_{k}\odot (g_{k}-\grad f(w_{k}))}^{2}\mid\mathcal{F}_{k}\right].
\tag{3}
\]

\item \textbf{Bound the variance term.}\\
Since $\eta_{k}\le \frac{\alpha}{\varepsilon}\mathbf{1}$ elementwise and $\|g_{k}-\grad f(w_{k})\|\le\sigma$ a.s.,
\[
\E\!\left[\norm{\eta_{k}\odot (g_{k}-\grad f(w_{k}))}^{2}\mid\mathcal{F}_{k}\right]
\le \frac{\alpha^{2}}{\varepsilon^{2}}\sigma^{2}d. \tag{4}
\]

\item \textbf{Combine (1)–(4).}\\
Taking expectation of (1) and using (2)–(4) yields
\[
\E\!\big[f(w_{k+1})\mid\mathcal{F}_{k}\big]
\le f(w_{k})-\big\langle\grad f(w_{k}),\eta_{k}\odot\grad f(w_{k})\big\rangle
+\frac{L}{2}\norm{\eta_{k}\odot\grad f(w_{k})}^{2}
+\frac{L\alpha^{2}\sigma^{2}d}{2\varepsilon^{2}}.
\tag{5}
\]

\item \textbf{Rewrite the inner product term.}\\
Because the division is elementwise,
\[
\big\langle\grad f(w_{k}),\eta_{k}\odot\grad f(w_{k})\big\rangle
= \sum_{j=1}^{d}\frac{\alpha\,\big(\grad_{j}f(w_{k})\big)^{2}}{\sqrt{v_{k+1,j}}+\varepsilon}
\ge
\frac{\alpha}{\sqrt{v_{k+1}^{\max}}+\varepsilon}\,\norm{\grad f(w_{k})}^{2},
\tag{6}
\]
where $v_{k+1}^{\max}:=\max_{j}v_{k+1,j}$.

\item \textbf{Upper bound the quadratic term.}\\
Similarly,
\[
\norm{\eta_{k}\odot\grad f(w_{k})}^{2}
= \sum_{j=1}^{d}\frac{\alpha^{2}\big(\grad_{j}f(w_{k})\big)^{2}}{(\sqrt{v_{k+1,j}}+\varepsilon)^{2}}
\le \frac{\alpha^{2}}{\varepsilon^{2}}\norm{\grad f(w_{k})}^{2}. \tag{7}
\]

\item \textbf{Plug (6)–(7) into (5).}\\
\[
\E\!\big[f(w_{k+1})\mid\mathcal{F}_{k}\big]
\le f(w_{k})
-\frac{\alpha}{\sqrt{v_{k+1}^{\max}}+\varepsilon}\norm{\grad f(w_{k})}^{2}
+\frac{L\alpha^{2}}{2\varepsilon^{2}}\norm{\grad f(w_{k})}^{2}
+\frac{L\alpha^{2}\sigma^{2}d}{2\varepsilon^{2}}.
\tag{8}
\]

\item \textbf{Choose $\alpha$ small enough.}\\
Assume $\alpha\le \frac{\varepsilon}{L}$, then $\frac{L\alpha^{2}}{2\varepsilon^{2}}\le\frac{\alpha}{2\varepsilon}$. Hence
\[
\frac{\alpha}{\sqrt{v_{k+1}^{\max}}+\varepsilon}
-\frac{L\alpha^{2}}{2\varepsilon^{2}}
\ge
\frac{\alpha}{2(\sqrt{v_{k+1}^{\max}}+\varepsilon)}.
\tag{9}
\]

\item \textbf{Simplify (8) using (9).}\\
\[
\E\!\big[f(w_{k+1})\mid\mathcal{F}_{k}\big]
\le f(w_{k})
-\frac{\alpha}{2(\sqrt{v_{k+1}^{\max}}+\varepsilon)}\norm{\grad f(w_{k})}^{2}
+\frac{L\alpha^{2}\sigma^{2}d}{2\varepsilon^{2}}.
\tag{10}
\]

\item \textbf{Take total expectation and telescope.}\\
Define $\Delta_{k}:=\E[f(w_{k})]$. Taking full expectation of (10) and summing $k=0$ to $T-1$ gives
\[
\Delta_{T}\le \Delta_{0}
-\frac{\alpha}{2}\sum_{k=0}^{T-1}\E\!\left[\frac{\norm{\grad f(w_{k})}^{2}}{\sqrt{v_{k+1}^{\max}}+\varepsilon}\right]
+\frac{L\alpha^{2}\sigma^{2}d}{2\varepsilon^{2}}T.
\tag{11}
\]

\item \textbf{Lower bound the denominator via Lemma~\ref{lem:sum_eta}.}\\
Since $v_{k+1}^{\max}\le\sum_{i=1}^{k+1}\|g_{i}\|^{2}\le G^{2}(k+1)$ by boundedness,
\[
\sqrt{v_{k+1}^{\max}}+\varepsilon\le G\sqrt{k+1}+\varepsilon\le 2G\sqrt{k+1}
\quad\text{for }k\ge1,
\]
so for all $k$,
\[
\frac{1}{\sqrt{v_{k+1}^{\max}}+\varepsilon}\ge\frac{1}{2G\sqrt{k+1}}.
\tag{12}
\]

\item \textbf{Insert (12) into (11).}\\
\[
\Delta_{T}\le \Delta_{0}
-\frac{\alpha}{4G}\sum_{k=0}^{T-1}\frac{\E[\norm{\grad f(w_{k})}^{2}]}{\sqrt{k+1}}
+\frac{L\alpha^{2}\sigma^{2}d}{2\varepsilon^{2}}T.
\tag{13}
\]

\item \textbf{Upper bound the harmonic sum.}\\
Using $\sum_{k=0}^{T-1}\frac{1}{\sqrt{k+1}}\ge 2\sqrt{T}$,
\[
\sum_{k=0}^{T-1}\frac{\E[\norm{\grad f(w_{k})}^{2}]}{\sqrt{k+1}}
\ge
\frac{1}{\sqrt{T}}\sum_{k=0}^{T-1}\E[\norm{\grad f(w_{k})}^{2}].
\tag{14}
\]

\item \textbf{Combine (13)–(14) and rearrange.}\\
\[
\frac{1}{T}\sum_{k=0}^{T-1}\E[\norm{\grad f(w_{k})}^{2}]
\le
\frac{4G\big(\Delta_{0}-\Delta_{T}\big)}{\alpha\sqrt{T}}
+\frac{2L\alpha\sigma^{2}d}{\varepsilon^{2}\sqrt{T}}.
\tag{15}
\]
Since $f(w_{T})\ge f^{\star}$, $\Delta_{0}-\Delta_{T}\le f(w_{0})-f^{\star}$.
Finally,
\[
\frac{1}{T}\sum_{k=0}^{T-1}\E[\norm{\grad f(w_{k})}^{2}]
\le
\underbrace{\frac{4G\big(f(w_{0})-f^{\star}\big)}{\alpha}}_{\displaystyle C_{1}}
\frac{1}{\sqrt{T}}
\;+\;
\underbrace{\frac{2L\alpha\sigma^{2}d}{\varepsilon^{2}}}_{\displaystyle C_{2}}
\frac{1}{\sqrt{T}}.
\tag{16}
\]

\item \textbf{Optimize the constant $\alpha$.}\\
Choosing $\alpha=\Theta\!\big(\varepsilon/\sqrt{L}\big)$ balances the two terms, yielding the bound stated in Theorem~\ref{thm:main}:
\[
\frac{1}{T}\sum_{k=0}^{T-1}\E[\norm{\grad f(w_{k})}^{2}]
\le
\frac{2L\big(f(w_{0})-f^{\star}\big)}{\alpha\sqrt{T}}
+\frac{2\alpha\sigma^{2}\sqrt{d}}{\sqrt{T}}.
\]
The additional $2\alpha G^{2}\sqrt{d}/\sqrt{T}$ term appears when we keep the deterministic bound $G$ instead of the variance bound.
\end{enumerate}

Thus the theorem holds. $\square$

\vspace{0.5em}
\section*{Rate Derivation (With Constants)}
From (16) we identify
\[
C_{1}= \frac{4G\big(f(w_{0})-f^{\star}\big)}{\alpha},
\qquad
C_{2}= \frac{2L\alpha\sigma^{2}d}{\varepsilon^{2}}.
\]
Setting $\alpha=\varepsilon/\sqrt{L}$ gives
\[
C_{1}=4G\sqrt{L}\frac{f(w_{0})-f^{\star}}{\varepsilon},
\qquad
C_{2}=2\varepsilon\sigma^{2}d\sqrt{L}.
\]
Both scale as $O(1/\sqrt{T})$, confirming the $O(\log T/\sqrt{T})$ bound after converting the average bound to a minimum‑gradient bound via Jensen’s inequality.

\vspace{0.5em}
\section*{Special Cases}
\begin{enumerate}
\item \textbf{Deterministic gradients ($\sigma=0$).} The variance term disappears and we obtain a pure $O(1/\sqrt{T})$ decay.
\item \textbf{Convex $f$.} Using convexity instead of the non‑convex gradient norm argument yields the classic AdaGrad regret bound $O\big(\frac{1}{\sqrt{T}}\sum_{j=1}^{d}\|w_{0,j}\|\big)$.
\item \textbf{Sparse gradients.} If only $s\ll d$ coordinates are non‑zero at each iteration, then $v_{k}^{\max}$ grows proportionally to $s$, effectively replacing $\sqrt{d}$ by $\sqrt{s}$ in the constants.
\end{enumerate}

\end{document}